\begin{question}

  \begin{questionpart}
    Let $x$ and $p_x$ be the coordinate and linear momentum in one
    dimension.  Evaluate the classical Poisson bracket
    \begin{equation*}
      \left[
        x,F(p_x)
        \right]_{\text{classical}}
    \end{equation*}
  \end{questionpart}

  \begin{questionpart}
    Let $x$ and $p_x$ be the corresponding quantum-mechanical
    operators this time.  Evaluate the commutator,
    \begin{equation*}
      \left[
        x,\exp(\frac{ip_xa}{\hbar})
        \right]
    \end{equation*}
  \end{questionpart}

  \begin{questionpart}
    Using the result obtained in (b), prove that
    \begin{equation*}
      \exp(\frac{ip_xa}{\hbar})\ket{x'} (x\ket{x'} = x'\ket{x'})
    \end{equation*}
    is an eigennstate of the coordinate operator x.  What is the
    corresponding eigenvalue?
  \end{questionpart}
\end{question}

\begin{purpose}
  We're going to dig deep in the consequences of the jog operator here
  and try to learn about about how it behaves when employed.
\end{purpose}

\begin{answer}

  \begin{answerpart}{Classical}
    Pulling in the definition of Poisson brackets from the book
    
    \begin{equation}
      \tag{1.230}
      [A(q,p),B(q,p)]_{\text{classical}} \equiv
      \sum_s \left(
      \pdv{A}{q_s}
      \pdv{B}{p_s}
      - 
      \pdv{A}{p_s}
      \pdv{B}{q_s}
      \right)
    \end{equation}

    $s$ appears to be the dimensionality which is specified to be $1$
    in the problem so we can rewrite that for our particular problem
    as follows

    \begin{equation}
      \begin{split}
        [x,F(p)]_{\text{classical}}
        &\equiv
        \pdv{x}{x}
        \pdv{F(p)}{p}
        - 
        \underbrace{
          \pdv{x}{p}
          \pdv{F(p)}{x}
          }_{0}
        \\
        &= \pdv{F(p)}{p}
      \end{split}
    \end{equation}

    Note that you can also make that negative if you want just by
    changing the order of the parameters even if you keep the order of
    the operators the same.  Hamiltonians require this order to get
    the right result without scratching your head.
  \end{answerpart}

  \begin{answerpart}{Quantum}

    There are three ways to solve this problem.  The first is, by far,
    the easisest which is to just employ Dirac's trick listed in the book

    \begin{equation}
      \tag{1.229}
      [\ ,\ ]_{\text{classical}} \rightarrow \frac{[\ ,\ ]}{\hbar}
    \end{equation}

    But...that's no fun...  Let's recall our exponentiation definition
    and dig in
 
    \begin{equation}
      e^{A} \equiv \sum_{k=0}^{\infty} \frac{1}{k!} A^k
    \end{equation}

    Apply to the problem
    
    \begin{equation}
      \begin{split}
        \left[x, \exp(\frac{i\hat{p}_xa}{\hbar})\right]
        &=
        x\exp(\frac{i\hat{p}_xa}{\hbar})
        -
        \exp(\frac{i\hat{p}_xa}{\hbar})x
        \\
        &=
          \left[x,
            \sum_{k=0}^\infty\frac{1}{k!}
            \left(\frac{i}{\hbar}a\hat{p}\right)^k
            \right]
        \\
        &=
        x
        \sum_{k=0}^\infty\frac{1}{k!}
        \left(\frac{i}{\hbar}a\hat{p}\right)^k
        -
        \sum_{k=0}^\infty\frac{1}{k!}
        \left(\frac{i}{\hbar}a\hat{p}\right)^k
        x
        \\
        &=
        \sum_{k=0}^\infty
        \frac{1}{k!}
        a^k
        \left(\frac{i}{\hbar}\right)^k
        \left(
        x\hat{p}^k
        -
        \hat{p}^kx
        \right)
        \\
      \end{split}
    \end{equation}

    \begin{equation}
      \begin{split}
        \left[x, \exp(\frac{i\hat{p}_xa}{\hbar})\right]
        &=
          \left[x,
            \sum_{k=0}^\infty\frac{1}{k!}
            \left(\frac{i}{\hbar}a\hat{p}_x\right)^k
            \right]
        \\
        &=
        \sum_{k=0}^{\infty}\frac{1}{k!}
        \left[
          x,\ 
          \left(\frac{i}{\hbar}a\hat{p}_x\right)^k
          \right]
        \\
        &=
        \sum_{k=0}^{\infty}\frac{1}{k!}
        \left(\frac{ia}{\hbar}\right)^k
        \left[
          x,\ 
          \hat{p}_x^k
          \right]
        \\
      \end{split}
    \end{equation}

    We then pull out and recursively apply a commutation relation from
    the book

    \begin{equation}
      \tag{1.232e}
      [A,BC] = [A,B]C + B[A,C]
    \end{equation}

    Which for us, becomes
    
    \begin{equation}
      [x,\hat{p}_x\hat{p}_x] = [x,\hat{p}_x]\hat{p}_x + \hat{p}_x[x,\hat{p}_x]
    \end{equation}

    Applied for powers of $n$

    \begin{equation}
      \inlabel{IHaveThePower}
      \begin{split}
        [x,\hat{p}_x^4]
        &= [x,\hat{p}_x\hat{p}_x^{3}]\\
        &= [x,\hat{p}_x]\hat{p}_x^{3} + \hat{p}_x[x,\hat{p}_x^{3}]\\
        &= [x,\hat{p}_x]\hat{p}_x^{3} + \hat{p}_x([x,\hat{p}_x]\hat{p}_x^{2} + \hat{p}_x[x,\hat{p}_x^{2}])\\
        &= [x,\hat{p}_x]\hat{p}_x^{3} + \hat{p}_x[x,\hat{p}_x]\hat{p}_x^{2} + \hat{p}_x^2[x,\hat{p}_x^{2}]\\
        &= [x,\hat{p}_x]\hat{p}_x^{3} + \hat{p}_x[x,\hat{p}_x]\hat{p}_x^{2} + \hat{p}_x^2([x,\hat{p}_x]\hat{p}_x + \hat{p}_x[x,\hat{p}_x])\\
        &= [x,\hat{p}_x]\hat{p}_x^{3} + \hat{p}_x[x,\hat{p}_x]\hat{p}_x^{2} + \hat{p}_x^2[x,\hat{p}_x]\hat{p}_x + \hat{p}_x^3[x,\hat{p}_x]\\
        &= \hat{p}_x^3[x,\hat{p}_x] + [x,\hat{p}_x]\hat{p}_x^{3} + \hat{p}_x[x,\hat{p}_x]\hat{p}_x^{2} + \hat{p}_x^2[x,\hat{p}_x]\hat{p}_x\\
        &= \hat{p}_x^{n-1}[x,\hat{p}_x] + \sum_{i=1}^{n-1}\hat{p}_x^{i-1}[x,\hat{p}_x]\hat{p}_x^{n-i}
      \end{split}
    \end{equation}

    Making use of the commutation of $[x,\hat{p}_x]$

    \begin{equation}
      \tag{1.228}
      [x_i,x_j] = 0,\quad[p_i,p_j] = 0,\quad[x_i,p_j] = i\hbar\delta_{ij}
    \end{equation}

    \begin{equation}
      \inlabel{ConstantPower}
      \begin{split}
        [x,\hat{p}_x^n]
        &= \hat{p}_x^{n-1}[x,\hat{p}_x] + \sum_{i=1}^{n-1}\hat{p}_x^{i-1}[x,\hat{p}_x]\hat{p}_x^{n-i}\\
        &= i\hbar \hat{p}_x^{n-1} + \sum_{i=1}^{n-1}i\hbar \hat{p}_x^{i-1}\hat{p}_x^{n-i}\\
        &= i\hbar
        \left(
        \hat{p}_x^{n-1} + \sum_{i=1}^{n-1} \hat{p}_x^{\cancel{i} - 1 + n - \cancel{i}}
        \right)
        \\
        &= i\hbar
        \left(
        \hat{p}_x^{n-1} + (n-1)\hat{p}_x^{n-1}
        \right)
        \\
        &= i \hbar n\hat{p}_x^{n-1}
      \end{split}
    \end{equation}

    Substituting that into our above expression

    \begin{equation}
      \inlabel{presub}
      \begin{split}
        \left[x, \exp(\frac{i\hat{p}_xa}{\hbar})\right]
        &=
        \sum_{k=0}^{\infty}\frac{1}{k!}
        \left(\frac{ia}{\hbar}\right)^k
        \left[
          x,\ 
          \hat{p}_x^k
          \right]
        \\
        &=
        \sum_{k=0}^{\infty}\frac{1}{k!}
        \left(\frac{ia}{\hbar}\right)^k
        \left(
        i \hbar k\hat{p}_x^{k-1}
        \right)
        \\
        &=
        i \hbar
        \sum_{k=0}^{\infty}
        \frac{k}{k!}
        \left(\frac{ia}{\hbar}\right)^k
        \left(
        \hat{p}_x^{k-1}
        \right)
        \\
        &=
        i \hbar
        \frac{ia}{\hbar}
        \underbrace{\sum_{k=1}^{\infty}\frac{1}{(k-1)!}}_{\text{0/0! = 0}}
        \left(\frac{ia}{\hbar}\hat{p}_x\right)^{k-1}
        \\
        &=
        -a
        \sum_{k=1}^{\infty}\frac{1}{(k-1)!}
        \left(\frac{ia}{\hbar}\hat{p}_x\right)^{k-1}
        \\
      \end{split}
    \end{equation}

    Doing a quick substitution for $k' = k-1$

    \begin{equation}
      \inlabel{postsub}
      \begin{split}
        \left[x, \exp(\frac{i\hat{p}_xa}{\hbar})\right]
        &=
        -a
        \sum_{k'=0}^{\infty}\frac{1}{k'!}
        \left(\frac{ia}{\hbar}\hat{p}_x\right)^{k'}
        \\
        &=
        -a
        \exp(\frac{ia}{\hbar}\hat{p}_x)
        \\
        &=
        \left(i\hbar\right)
        \left(
        \frac{ia}{\hbar}
        \exp(\frac{ia}{\hbar}\hat{p}_x)
        \right)
        \\
        &=
        \left(i\hbar\right)
        \left(
        \pdv{}{p}
        \exp(\frac{ia}{\hbar}\hat{p}_x)
        \right)
        \\
        &=
        i\hbar
        \pdv{}{p}
        \exp(\frac{ia}{\hbar}\hat{p}_x)
        \\
      \end{split}
    \end{equation}

    So, in the end

    \begin{equation}
      \inlabel{xfp}
      \begin{split}
        [x, F(\hat{p}_x)]
        &= \left[x, \exp(\frac{i\hat{p}_xa}{\hbar})\right]\\
        &= 
        i\hbar
        \pdv{}{p}
        \exp(\frac{ia}{\hbar}\hat{p}_x)
        \\
        &=
        i\hbar
        \pdv{}{p_x}
        F(\hat{p}_x)\qed
        \\
      \end{split}
    \end{equation}

    Pulling in the definition of the jog operator

    \begin{equation}
      \tag{1.218}
      \mathscr{J}(\Delta x' \vu{x}) = \exp(- \frac{ip_x\Delta x'}{\hbar})
    \end{equation}

    We see that this is exactly our result with $a=-\Delta x'$.
    Evaluating that derivative, we find that

    \begin{equation}
      \begin{split}
        [x, \mathscr{J}(\Delta x')]
        &=
        i\hbar
        \pdv{}{p_x}
        \mathscr{J}(\Delta x')
        \\
        &=
        \left( i\hbar \right)
        \left( \frac{-i\Delta x'}{\hbar} \right)
        \mathscr{J}(\Delta x')
        \\
        &=
        \Delta x'
        \mathscr{J}(\Delta x')
        \\
      \end{split}
    \end{equation}

    Now we take a look at a very similar equation in the book which is

    \begin{equation}
      \tag{1.207}
      [\vb{x}, \mathscr{J}(d\vb{x}')] = d\vb{x}'
    \end{equation}

    This was a bit alarming to me when I noticed the lack of jog
    operator on the right hand side, but $d\vb{x}'$ implies an
    infinitesimal jog which implies that the exponential is
    approximately $1$.  On the other hand, when we apply the operator
    at the macroscopic scale with $\Delta x'$, we get a result of
    scaling the operator by that same jog amount.  What does it mean
    to scale the operator like that?  If you look in section 1.6.4,
    you'll find a whole discussion about jogging $x$ then $y$ vs
    jogging $y$ then $x$ and how they're equivalent.  Given the
    exponential nature of the operator, it's pretty easy to see how a
    multiplicative factor outside becomes an additive factor
    inside\footnote{Technically a factor of $\ln(\Delta x').$} and how
    jogging up then right becomes equivalent to jogging right then up.

    There are other observations to be made from this result such as
    the consequences of observing then jogging vs jogging then
    observing.  If we take to heart this concept that a multiplicative
    factor outside the exponential is equivalent to an additive factor
    in the system, we see that

    \begin{equation}
      \begin{alignedat}{2}
        &&[x, \mathscr{J}(\Delta x')]
        &= x\mathscr{J}(\Delta x') - \mathscr{J}(\Delta x')x\\
        &&&=
        \Delta x'
        \mathscr{J}(\Delta x')
        \\
        \implies
        && x\mathscr{J}(\Delta x') - \mathscr{J}(\Delta x')x
        &= 
        \Delta x'
        \mathscr{J}(\Delta x')
        \\
      \end{alignedat}
    \end{equation}

    If we interpret that as ``The difference between observing then
    jogging vs jogging then observing is $\Delta x'$'' then it starts
    to make some amount of sense.
  \end{answerpart}

  \begin{answerpart}{All Together Now}
    Showing the eigen-nature of $\mathscr{J}(a)$ in the $x$ basis is now much
    simpler.

    \begin{equation}
      \begin{alignedat}{2}
        &&a\mathscr{J}(a)
        &= x\mathscr{J}(a) - \mathscr{J}(a)x
        \\
        \implies
        &&x\mathscr{J}(a)
        &= a\mathscr{J}(a) + \mathscr{J}(a)x
        \\
        \implies
        &&x\mathscr{J}(a)\ket{x'}
        &= (a\mathscr{J}(a) + \mathscr{J}(a)x)\ket{x'}
        \\
        &&&= (a + x')\mathscr{J}(a)\ket{x'}
        \\
      \end{alignedat}
    \end{equation}

    implying an eigenvalue of $a + x'$.
  \end{answerpart}

\end{answer}
